\chapter{APIs}
\label{ch:api}

\index{API}
\index{CSV file}
\index{database}

Back in Chapter~\ref{ch:screenScraping1}, I called the screen-scraping scenario
``the worst of all possible worlds.'' Using an \textbf{API}, the subject of
this lesson, is sort of the best. I say ``sort of'' because it's certainly not
the \textit{easiest}. The easiest is when you have a ready-made CSV file with
exactly the data you want, and you suck it into a \texttt{DataFrame}, done. The
thing is, much of the data you'll be interested in analyzing as a Data
Scientist wasn't put together like that for you. Instead, it exists as part of
a vast database somewhere, and you need flexibility and power to extract on
demand exactly the parts you want. This is what a good API lets you do.

API actually stands for ``Application Programming Interface,'' which is a
pretty uninformative acronym. What it really means, for our purposes, is ``a
documented way for one program to ask another program for data or services.''
As you'll see with experience, some APIs are better designed than others,
and make it easier to get what you want. In this chapter we'll work through one
main example -- NASA's live feed of near-Earth asteroids, no less -- and
generalize some common principles from it.

\section{REST APIs}
\index{REST API}

Another strange pseudo-acronym you'll encounter is
``\textbf{REST}''\footnote{This sort of stands for ``REpresentational State
Transfer,'' another utterly unhelpful phrase.}, which is a \textit{style} of
API. REST APIs have the following characteristics:

\begin{itemize}
\itemsep.1em

\index{URL}

\item They're designed to be used over the Web. (Although ``over the Web'' may
not mean what you think it means; for one, it \textit{doesn't} mean ``using a
browser.'' Instead, it means ``accessing the data using HTTP requests in the
form of structured URLs/hyperlinks.'' We'll be writing \textit{Python code} to
retrieve what appear to be pseudo-Web-pages outside our browser.)

\index{SOAP API}
\index{RPC}
\index{Fielding, Roy}

\item They're generally simple and lightweight. REST became popular in part
as an alternative to heavier Web-service approaches such as SOAP. It emphasizes
a uniform interface and the ordinary conventions of the Web, which tends to
make REST APIs pretty nice to use.

\index{HTTP protocol}

\item REST APIs on the Web are normally built directly on the \textbf{HTTP}
protocol, and make use of HTTP methods like \texttt{GET}, \texttt{HEAD},
\texttt{POST}, \texttt{DELETE}, \texttt{PUT}, \textit{etc.} Because of this,
different REST APIs tend to have a lot in common in how they work.

\index{stateless (protocol)}
\index{server}
\index{client}

\item They're ``\textbf{stateless},'' as HTTP is itself. Roughly speaking,
this means that each request has to stand on its own: the server shouldn't
need to remember some detail about an earlier request in order to understand
the current one. The server can certainly maintain databases and other
persistent data; it's the client's \textit{session state} that isn't carried
implicitly from one request to the next. (``\textbf{Server}'' is a word
meaning ``a machine on the Internet that you want to access,'' and
``\textbf{client}'' here essentially means ``your own computer, which makes
requests of the server.'')

\index{JSON}
\index{XML}
\item The data they provide can come in a variety of formats, typically JSON
or XML.

\end{itemize}

\subsection{Other terms}

\index{RESTful@``RESTful'' API}

APIs that follow the constraints of Roy Fielding's REST architectural style
are known colloquially as -- I kid you not -- ``\textbf{RESTful}.''

\index{Web services}
\index{SOA (Service-Oriented Architecture)}

You'll also sometimes hear the buzzwords ``\textbf{Web services}'' or
``\textbf{SOA}'' (Service-oriented Architecture). Although these terms are a
bit broader and communicate something subtly different than just ``the use of
Web-based APIs'' (REST or otherwise), in practice you will see terms like this
used jointly to refer to the same technologies. If you hear someone talking
about ``a Web services approach,'' your ears should perk up and you should
anticipate being able to apply at least some of the lessons in this chapter.

\index{endpoint}

Finally, the term \textbf{endpoint} refers to a particular, addressable part
of an API that is normally defined by a base URL. For instance, NASA's
near-Earth-object ``feed'' endpoint is at
\url{https://api.nasa.gov/neo/rest/v1/feed}, and the endpoint for looking up
everything about one particular asteroid is
\url{https://api.nasa.gov/neo/rest/v1/neo} followed by that asteroid's ID. In
general, an API will have many different endpoints, each of which defines its
own parameters and is used for a particular type of information request.

\section{RESTful APIs and the \texttt{requests} package}

\index{import@\texttt{import}}
\index{requests@\texttt{requests} package}

Python's \texttt{requests} package is an easy way to make a call to an
external API and get the results. It basically encapsulates the under-the-hood
stuff that formulates, sends, and parses the results of an HTTP request.

\index{get@\texttt{get()} function (\texttt{requests})}

HTTP defines various operations to communicate with a server such as
\texttt{GET}, \texttt{POST}, and \texttt{HEAD}. The \texttt{requests} package
bundles these up in functions like: \texttt{requests.get()},
\texttt{requests.post()}, \textit{etc.} The most common method we'll use in API
access is the \texttt{.get()} command. It can actually be used to retrieve
\textit{any} Web-accessible page on the network:

\begin{Verbatim}[fontsize=\small,samepage=true,frame=single,framesep=3mm]
import requests
class_home_page = requests.get(
    "http://stephendavies.org/data219")
print(class_home_page.text)
\end{Verbatim}

\index{HTML}
\index{text@\texttt{.text} (\texttt{requests} package)}
\index{Unicode}

That ``\texttt{.text}'' on the end gives you the actual contents of the
response -- the HTML code, in this case -- decoded into one big ordinary
Python string. Character encodings (like \textbf{Unicode}) are an interesting
subject in their own right, but happily, \texttt{requests} normally figures
out the appropriate encoding for text responses from the information the
server sends along with them.\footnote{If you want the response as raw bytes
instead -- say, because you're downloading an image -- ask for
\texttt{.content}. And when a response contains JSON, as most of the ones
later in this chapter will, we'll normally use \texttt{.json()} instead.}

Anyway, if you run the above code, you'll see the (gnarly) HTML code for the
entire course web page on your screen as one big string. Seeing this emphasizes
that browsers and REST APIs are using much of the same underlying machinery:
both send HTTP requests to servers and receive responses. A browser commonly
asks for HTML that it can render as a Web page; an API client commonly asks for
machine-readable data, such as JSON. This shared HTTP foundation is part of
what makes REST APIs so straightforward to work with.

\section{The NASA NeoWs API}
\index{NeoWs (Near Earth Object Web Service)}
\index{NASA}
\index{asteroid}

Most of the time we won't get HTML back, but rather a data format such as XML
or JSON. Our running example will be NASA's \textbf{NeoWs} (``Near Earth
Object Web Service''), a genuinely live database of the asteroids that pass
near our planet -- what they're named, how big they are, how close they'll
come, and whether NASA has flagged them as \textit{potentially hazardous}. I
submit that this beats any CSV file you've ever been handed.

\index{Apophis (asteroid)}

Let's start by asking it about a single asteroid: number 2099942, better known
as \textbf{Apophis}, briefly famous in 2004 for having a scarier-than-usual
chance of hitting us.\footnote{It won't. Astronomers have since ruled out an
impact for at least a century. You can relax and learn APIs.} 

\begin{Verbatim}[fontsize=\small,samepage=true,frame=single,framesep=3mm]
response = requests.get("https://api.nasa.gov/neo/rest/" +
    "v1/neo/2099942?api_key=DEMO_KEY")
print(response.text)
\end{Verbatim}
\vspace{-.2in}

(More on that ``\texttt{api\_key}'' business in a moment.)


The beginning of the response looks something like this (I've added line
breaks and omitted some fields so that it fits on the page):

\begin{Verbatim}[fontsize=\scriptsize,samepage=true,frame=leftline,framesep=5mm,framerule=1mm]
{"links":{"self":"http://api.nasa.gov/neo/rest/v1/neo/2099942?
api_key=DEMO_KEY"},
 "id":"2099942",
 "neo_reference_id":"2099942",
 "name":"99942 Apophis (2004 MN4)",
 "designation":"99942",
 "nasa_jpl_url":"https://ssd.jpl.nasa.gov/tools/sbdb_lookup.html
#/?sstr=2099942",
 ...
 "is_potentially_hazardous_asteroid":true,
 "close_approach_data":[ ... ]}
\end{Verbatim}


\index{pretty printing@``pretty printing''}
\index{pprint()@\texttt{pprint()} function}
\index{json()@\texttt{.json()} (\texttt{requests} package)}

That's the real thing, but it's awfully hard on the eyes as one giant string.
Before we use the result, \texttt{raise\_for\_status()} checks whether the
server reported an HTTP error. Since the response is JSON, we can then ask
\texttt{requests} to convert it to ordinary Python objects with the
\texttt{.json()} method,\footnote{\texttt{.json()} is essentially a
convenience wrapper around JSON decoding; here the result happens to be a
dictionary.} and pretty-print it:

\begin{Verbatim}[fontsize=\small,samepage=true,frame=single,framesep=3mm]
from pprint import pprint
response.raise_for_status()
apophis = response.json()
pprint(apophis)
\end{Verbatim}

\begin{Verbatim}[fontsize=\small,samepage=true,frame=leftline,framesep=5mm,framerule=1mm]
{'absolute_magnitude_h': 19.09,
 'close_approach_data': [ ...loads of stuff... ],
 'designation': '99942',
 'estimated_diameter': {'kilometers':
      {'estimated_diameter_max': 0.6134,
       'estimated_diameter_min': 0.2743},
      ...more units... },
 'id': '2099942',
 'is_potentially_hazardous_asteroid': True,
 'links': {'self': 'http://api.nasa.gov/neo/rest/v1/neo/
2099942?api_key=DEMO_KEY'},
 'name': '99942 Apophis (2004 MN4)',
 'nasa_jpl_url': 'https://ssd.jpl.nasa.gov/tools/
sbdb_lookup.html#/?sstr=2099942',
 'neo_reference_id': '2099942',
 'orbital_data': { ...loads more stuff... }}
\end{Verbatim}

\index{dictionary}

Better. (I've elided a couple of the giant inner structures so this fits on a
page.) You can see that the API gave us a well-structured chunk of data, in a
dictionary, which we can explore and manipulate in all the usual ways --
including the nested drilling-down we practiced in the JSON chapters. 
Do notice that the \\
\texttt{is\_potentially\_hazardous\_asteroid} field is \texttt{True}.
I told you this was more exciting than a CSV file.

\index{URL}

Notice also that some of the data consists of URLs. The \texttt{links} entry
points back to this API resource, while \texttt{nasa\_jpl\_url} points instead
to a human-readable JPL Web page with further information about Apophis. These
aren't quite the same thing: one is meant for programs like ours, while the
other is mainly meant for people. We'll see in a moment that NeoWs also embeds
API URLs that let us \textit{page} through large result sets.

\section{API keys and rate limits}

\subsection{Authentication}

\index{authentication}
\index{API key}
\index{rate limit}

Many APIs require some form of \textbf{authentication}, which means that your
request includes some credential identifying you, your application, or both.
With an open-data API like NASA's, the key isn't there to protect secret
asteroid data from you; among other things, it lets the service identify API
usage and enforce \textbf{rate limits} that constrain how many requests one
user can make in a given amount of time.

Every site has a different procedure for this, and the documentation varies
considerably in helpfulness. NASA's, mercifully, is about the simplest one in
existence:

\begin{compactenum}
\item Go to \url{https://api.nasa.gov}.
\item Type in your name and email address.
\item Your very own \textbf{API key} -- a string of gobbledygook like
\texttt{6ZbY6...} (mine goes on for forty characters) -- arrives in your inbox
moments later.
\end{compactenum}

\index{DEMO\_KEY@\texttt{DEMO\_KEY}}

You'll notice our Apophis request didn't even need your own key: NASA provides
a communal training-wheels key, literally called \texttt{DEMO\_KEY}, that
anyone can use without signing up. As of this writing, it's limited to 30
requests per IP address per hour and 50 per day. That's fine for kicking the
tires, but it's easy to exhaust -- especially in a classroom, where many
computers may appear to NASA under the same public IP address. A registered
NASA API key normally gets a much larger allowance, so get your own if you're
going to do anything substantial.

Unlike fancier schemes, this style of authentication couldn't be simpler to
\textit{use}: you just pass your key as one more parameter (called
\texttt{api\_key}) in each request, exactly as we did above. Your personal API
key should nevertheless be treated like a password: don't post it publicly or
commit it to a shared Git repository. We'll continue using
\texttt{DEMO\_KEY} in the printed examples so that no real credential ever
appears in the book.

\subsection{Watching your rate limit}

\index{headers (HTTP)}

How do you know where you stand with your hourly allowance? The server tells
you, in the \textbf{headers} it attaches to every response -- a little
envelope of metadata that rides along with the actual data:

\begin{Verbatim}[fontsize=\small,samepage=true,frame=single,framesep=3mm]
print(response.headers['X-RateLimit-Limit'])
print(response.headers['X-RateLimit-Remaining'])
\end{Verbatim}
\vspace{-.2in}

\begin{Verbatim}[fontsize=\small,samepage=true,frame=leftline,framesep=5mm,framerule=1mm]
30
27
\end{Verbatim}

\index{caching}

Thirty an hour on \texttt{DEMO\_KEY}, and I've already burned three of mine
writing this chapter. The professional habit that follows from this: make your
API call \textit{once}, save the JSON you get back to a file, and develop your
analysis against your saved copy rather than hammering the server on every
run of your program. This is kinder to NASA, kinder to your rate limit, and --
bonus -- your code still works when the WiFi doesn't.

One wrinkle to watch for: NeoWs puts the API key back into some of the URLs
inside its JSON response. If you're using your personal key, therefore, don't
post the raw saved response publicly without checking it first.

\index{OAuth}

One last word on authentication and authorization: some APIs -- especially
ones involving a user's private data -- use a fancier, multi-step system
called \textbf{OAuth} (see \url{oauth.net}). OAuth is an open standard for
\textit{authorization}: it lets you grant an application limited access to
data or services on your behalf without simply handing that application your
password. (You've probably used it without realizing it, whenever some
third-party app popped up a dialog saying ``application so-and-so is
requesting access to the following information of yours...do you approve?'')
The good news is that when you eventually need OAuth, an API's official client
library or another specialized package will often handle much of the ceremony
for you.

\section{HTTP GET parameters}
\index{HTTP}
\index{get@\texttt{GET} parameter (HTTP)}

The NeoWs API is documented at \url{https://api.nasa.gov}, a page well worth
perusing in detail. You'll see its endpoints documented there, each of which
has its parameters listed.

Conceptually, making an API request is like making a function call: you're
triggering some operation, passing parameters, and getting a return value. The
syntax is a bit different though. In Python, you're used to making a function
call like this:

\begin{Verbatim}[fontsize=\scriptsize,samepage=true,frame=single,framesep=3mm]
get_asteroids('2026-03-01','2026-03-07')<--(not real; I'm making this up)
\end{Verbatim}

to indicate something like ``get all the asteroids passing near Earth during
the first week of March.'' Now in order to ``call'' this ``function'' from the
\textit{API} and ``pass'' these ``parameters,'' you code up a URL with syntax
like this:

\begin{center}
\scriptsize{(all one string:)}
\end{center}
\normalsize
\vspace{-.5in}
\begin{Verbatim}[fontsize=\small,samepage=true,frame=lines,framesep=3mm]
https://api.nasa.gov/neo/rest/v1/feed?start_date=2026-03-01&
end_date=2026-03-07&api_key=DEMO_KEY
\end{Verbatim}

Stare at all that nonsense. Note that our ``parameters'' are expressed as
key/value pairs, with no spaces, at the end of the URL: after a ``\texttt{?}'',
and each separated by ``\texttt{\&}'' signs.

\index{params@\texttt{params} argument (\texttt{requests})}

Now, you \textit{can} build such URLs by gluing strings together, and it's
worth doing a couple times just to feel the machinery. But the
\texttt{requests} package will happily do the punctuation for you: hand
\texttt{.get()} a plain dictionary of your parameters via its \texttt{params}
argument, and it assembles the URL itself:

\begin{Verbatim}[fontsize=\small,samepage=true,frame=single,framesep=3mm]
response = requests.get(
    "https://api.nasa.gov/neo/rest/v1/feed",
    params={ 'start_date': '2026-03-01',
             'end_date': '2026-03-07',
             'api_key': 'DEMO_KEY' },
    timeout=10)
response.raise_for_status()
print(response.url)
feed = response.json()
\end{Verbatim}

The \texttt{timeout=10} is another good habit: it tells
\texttt{requests} not to sit around forever if the server stops responding.

\begin{Verbatim}[fontsize=\scriptsize,samepage=true,frame=leftline,framesep=5mm,framerule=1mm]
https://api.nasa.gov/neo/rest/v1/feed?start_date=2026-03-01&
end_date=2026-03-07&api_key=DEMO_KEY
\end{Verbatim}

Same URL, zero chances to mistype an ampersand. Let's see what a week of
asteroids looks like, structurally:

\begin{Verbatim}[fontsize=\small,samepage=true,frame=single,framesep=3mm]
print(feed['element_count'])
print(list(feed['near_earth_objects'].keys()))
\end{Verbatim}
\vspace{-.2in}

At the time this example was prepared, the result looked like this:

\begin{Verbatim}[fontsize=\small,samepage=true,frame=leftline,framesep=5mm,framerule=1mm]
40
['2026-03-01', '2026-03-02', '2026-03-03', '2026-03-04',
 '2026-03-05', '2026-03-06', '2026-03-07']
\end{Verbatim}

Forty asteroids in one week, and note the API designers' interesting
structural choice: \texttt{near\_earth\_objects} is a dictionary whose
\textit{keys are the dates}, each mapped to a list of that day's asteroids.
File that fact away; we'll need it when we flatten this thing at the end of
the chapter.

(A heads-up: when \textit{you} run this, don't expect exactly forty, even
though these dates are fixed in the past. The catalog behind this API is
alive: small asteroids are routinely discovered \textit{after} they've
already flown past us, and each new discovery is filed retroactively into
the date of its approach, so a ``historical'' week quietly grows over time.
Your query is frozen; the database isn't. Which is one more argument for
the habit from earlier in this chapter -- save the response to a file, and
your analysis stands on data that holds still.)

\section{URL encoding}

Sometimes the parameters you want to put into a URL contain spaces,
punctuation, or other characters that can't simply be dropped into a query
string as-is. These values have to be \textbf{URL encoded}. Depending on the
encoding being used, for instance, a space may appear as ``\texttt{+}'' or as
``\texttt{\%20}'', and other special characters are commonly represented by
a percent sign followed by two hexadecimal digits.

\index{percent syntax}
\index{urllib@\texttt{urllib} package}
\index{quote\_plus@\texttt{quote\_plus()} function (\texttt{urllib})}

You normally won't need to do this by hand. Python's
\texttt{quote\_plus()} function illustrates the idea:

\begin{Verbatim}[fontsize=\footnotesize,samepage=true,frame=single,framesep=3mm]
import urllib.parse
print(urllib.parse.quote_plus(
    "star wars toys! 20% off"))
\end{Verbatim}
\vspace{-.2in}

\begin{Verbatim}[fontsize=\small,samepage=true,frame=leftline,framesep=5mm,framerule=1mm]
star+wars+toys%21+20%25+off
\end{Verbatim}

Here \texttt{quote\_plus()} encoded one \textit{parameter value}; we wouldn't
run an entire API URL through it. And in ordinary \texttt{requests} code,
there's an even easier answer: pass your parameter values through
\texttt{params=} as above and let \texttt{requests} encode them for you.
That's what it's for.

\section{Paging through results}
\index{pagination}

A courteous API won't dump ten years of data on you in one response, no matter
how nicely you ask. Instead it caps each response at a sane size and gives you
a way to ask for the \textit{next} chunk. NeoWs, for instance, will hand over
at most one week of asteroids per request -- and look what else was hiding in
our feed:

\begin{Verbatim}[fontsize=\small,samepage=true,frame=single,framesep=3mm]
pprint(feed['links'])
\end{Verbatim}
\vspace{-.2in}

\begin{Verbatim}[fontsize=\scriptsize,samepage=true,frame=leftline,framesep=5mm,framerule=1mm]
{'next': 'http://api.nasa.gov/neo/rest/v1/feed?start_date=
2026-03-07&end_date=2026-03-13&api_key=DEMO_KEY',
 'prev': 'http://api.nasa.gov/neo/rest/v1/feed?start_date=
2026-02-23&end_date=2026-03-01&api_key=DEMO_KEY',
 'self': 'http://api.nasa.gov/neo/rest/v1/feed?start_date=
2026-03-01&end_date=2026-03-07&api_key=DEMO_KEY'}
\end{Verbatim}

There are those embedded REST URLs again: the API has pre-built the requests
for the \textit{following} week and the \textit{previous} week and handed them
to us, ready to use. Collecting a longer stretch of data is therefore just a
loop -- follow \texttt{next} until you've had enough:

\begin{Verbatim}[fontsize=\small,samepage=true,frame=single,framesep=3mm]
url = "https://api.nasa.gov/neo/rest/v1/feed?" + \
    "start_date=2026-03-01&api_key=DEMO_KEY"
weeks = []
for i in range(4):                    # four weeks' worth
    response = requests.get(url, timeout=10)
    response.raise_for_status()
    page = response.json()
    weeks.append(page)
    url = page['links']['next'].replace(
        "http://", "https://", 1)
\end{Verbatim}

That last \texttt{replace()} is a NeoWs-specific wrinkle: at the time of this
writing, the service puts \texttt{http://} rather than \texttt{https://} in
its embedded paging links. We'll switch it back to HTTPS before making the
next request.

Other APIs do the same job with \texttt{limit} and \texttt{offset} parameters
instead (``give me 50 results, starting at result 200''), but the idea is
identical: keep asking, politely, until you have what you need. Do remember
that every page is a fresh request as far as your rate limit is concerned --
which is one more reason to fetch once and save your results to a file.

\section{From JSON to \texttt{DataFrame}}
\index{DataFrame@\texttt{DataFrame} (Pandas)}
\index{flattening (JSON)}

Time for the payoff. A nested dictionary is nice to poke around in, but our
analysis tools want a good old flat \texttt{DataFrame}: one row per asteroid,
one column per attribute. The date-keyed structure we noted earlier tells us
exactly what to write -- an outer loop over the dates, an inner loop over each
date's asteroids:

\begin{Verbatim}[fontsize=\small,samepage=true,frame=single,framesep=3mm]
import pandas as pd
rows = []
for date, asts in feed['near_earth_objects'].items():
    for a in asts:
        rows.append({
            'date': date,
            'name': a['name'],
            'hazardous':
                a['is_potentially_hazardous_asteroid'],
            'miss_km': float(a['close_approach_data'][0]
                ['miss_distance']['kilometers']),
        })
asteroids = pd.DataFrame(rows)
print(asteroids.sort_values('miss_km').head(3))
\end{Verbatim}
\vspace{-.2in}

For this particular snapshot of the data, the three smallest distances were:

\begin{Verbatim}[fontsize=\small,samepage=true,frame=leftline,framesep=5mm,framerule=1mm]
          date            name  hazardous       miss_km
41  2026-03-05    (2026 DQ2)        False  8.913411e+05
7   2026-03-01    (2019 XQ1)        False  2.317744e+06
55  2026-03-07    (2013 FA7)         True  4.960023e+06
\end{Verbatim}

\index{float()@\texttt{float()} function}

One thing deserves a callout: that \texttt{float()}. NeoWs sends its
miss-distances as \textit{strings} -- \texttt{"891341.1"}, quotes and all --
so before you can sort or do arithmetic on them, you must convert them
yourself. Welcome to real data. One \texttt{float()} and we're respectable
again.

And that's the whole trip: from a live NASA service on the open Internet to a
sorted \texttt{DataFrame} of that week's closest approaches, in about a dozen
lines of Python. This is the reason APIs earn their ``sort of the best''
title.

\section{Specialized packages}

For many popular APIs, you may find Python packages written to streamline and
hide some of the protocol idiosyncrasies -- wrappers exist for everything from
Reddit to Spotify to, yes, NASA itself. These may be easier to use, depending
on how well written and maintained they are. Underneath, though, wrappers for
REST APIs are generally handling the same machinery you now understand:
endpoints, parameters, authentication, HTTP requests, and structured responses
such as JSON. That so many such packages exist is a testimony to how powerful
(and common) it is to be able to access data through external APIs.
